\section*{\texorpdfstring{\(\mathfrak{sl}_2\) -TRIPLETS}{\textbackslash mathfrak\{sl\}\_2 -TRIPLETS}}

\begin{enumerate}
\def\labelenumi{\arabic{enumi})}
\setcounter{enumi}{32}
\item
  An \(sl_{2}\) -triplet in g is a sequence \((x,h,y)\) of elements of g distinct from \((0,0,0)\) and such that \([h,x]=2x,[h,y]=-2y,[x,y]=-h\) . Then x,y are nilpotent in g, and h is semi-simple in g.
\item
  Let \(x\) be a non-zero nilpotent element of \(\mathfrak{g}\) . There exist \(h, y \in \mathfrak{g}\) such that \((x, h, y)\) is an \(\mathfrak{sl}_2\) -triplet.
\item
  Let \((x,h,y)\) and \((x',h',y')\) be \(\mathfrak{sl}_2\) -triplets in \(\mathfrak{g}\) . The following conditions are equivalent:
\end{enumerate}

\begin{enumerate}
\def\labelenumi{\alph{enumi})}
\item
  there exists \(s \in \text{Aut}_{e}(\mathfrak{g})\) such that \(sx = x'\) ;
\item
  there exists \(s \in \mathrm{Aut}_e(\mathfrak{g})\) such that \(sx = x', sh = h', sy = y'\) .
\end{enumerate}

\begin{enumerate}
\def\labelenumi{\arabic{enumi})}
\setcounter{enumi}{35}
\tightlist
\item
  If \(k\) is algebraically closed, conditions \(a)\) and \(b)\) of 35) are equivalent to: \(c)\) there exists \(s \in \operatorname{Aut}_e(\mathfrak{g})\) such that \(sh = h'\) .
\end{enumerate}

Moreover, the number of conjugacy classes, relative to \(\operatorname{Aut}_{e}(\mathfrak{g})\) , of non-zero nilpotent elements of g is at most equal to \(3^{l}\) , where l is the rank of g.

\begin{enumerate}
\def\labelenumi{\arabic{enumi})}
\setcounter{enumi}{36}
\tightlist
\item
  A nilpotent element x of g is called principal if the dimension of the centralizer of x is equal to the rank of g. There exist principal nilpotent elements in g. If k is algebraically closed, the principal nilpotent elements of g are conjugate under \(\operatorname{Aut}_{e}(\mathfrak{g})\) .
\end{enumerate}

